\documentclass[11pt]{article}
\usepackage[margin=1in]{geometry}
\usepackage{amsmath, amssymb, amsthm, bm, mathtools}
\usepackage{booktabs}
\usepackage{enumitem}
\usepackage{microtype}
\usepackage{hyperref}
\usepackage[most]{tcolorbox}

\hypersetup{colorlinks=true, linkcolor=blue, urlcolor=blue, citecolor=blue}

% ---------- handy macros ----------
\newcommand{\R}{\mathbb{R}}
\newcommand{\dd}{\,\mathrm{d}}
\newcommand{\ve}[1]{\mathbf{#1}}
\newcommand{\pd}[2]{\frac{\partial #1}{\partial #2}}
\newcommand{\dtotal}[2]{\frac{\mathrm{d} #1}{\mathrm{d} #2}}
\newcommand{\dqd}[2]{\frac{\mathrm{d} #1}{\mathrm{d} #2}}
\newcommand{\abs}[1]{\left\lvert #1 \right\rvert}
\newcommand{\paren}[1]{\left( #1 \right)}
\newcommand{\brak}[1]{\left[ #1 \right]}
\newcommand{\curly}[1]{\left\{ #1 \right\}}
\newcommand{\norm}[1]{\left\lVert #1 \right\rVert}
\DeclareMathOperator{\sgn}{sgn}

% ---------- theorem-like environments ----------
\newtheorem{definition}{Definition}
\newtheorem{proposition}{Proposition}
\newtheorem{example}{Example}
\newtheorem{remark}{Remark}

\title{Comparative Statics and the Concept of Derivative} %\\ \large After Chiang \& Wainwright (4th ed., Ch.~6) with Links to Standard Texts}
\author{ECON 320 - Introduction to Mathematical Economics}
\date{Fall 2025}

\tcbset{
  sectionbox/.style={
    colback=blue!10!white,
    colframe=blue!50!black,
    boxrule=0.8pt,
    arc=4pt,
    left=3pt,
    right=3pt,
    top=3pt,
    bottom=3pt
  },
  examplebox/.style={
    colback=green!10!white,
    colframe=green!50!black,
    boxrule=0.8pt,
    arc=4pt,
    left=3pt,
    right=3pt,
    top=3pt,
    bottom=3pt
  }
}
  
\begin{document}
\maketitle

\begin{tcolorbox}[sectionbox]
\textbf{Learning objectives.} By the end of this Lecture you should be able to:
\begin{enumerate}[label=\arabic*.]
  \item Define and interpret the derivative as a limit and as a marginal concept in economics.
  \item Apply rules of differentiation (product, quotient, chain) and compute higher-order derivatives.
  \item Use partial and total derivatives to analyze how endogenous variables respond to parametric changes.
  \item Perform comparative statics in single-equation models: \(F(x,p)=0 \Rightarrow \dtotal{x}{p} = -\dfrac{F_p}{F_x}\).
  \item Perform comparative statics in systems: \(\ve F(\ve x,\ve p)=\ve 0 \Rightarrow \dtotal{\ve x}{\ve p}=-J^{-1}B\) where \(J=\pd{\ve F}{\ve x}\), \(B=\pd{\ve F}{\ve p}\).
  \item Use Cramer's rule for sign analysis and interpret results economically.
  \item Work through canonical applications (supply--demand, tax incidence, monopoly, Cournot, Cobb--Douglas demand).
\end{enumerate}
\end{tcolorbox}

\newpage
\tableofcontents

\newpage

\section{Why derivatives and comparative statics matter}
Economic models map \emph{assumptions} (parameters, exogenous variables) into \emph{predictions} (endogenous variables). \textbf{Comparative statics} studies how equilibrium outcomes change when assumptions change. The \textbf{derivative} is the central tool: it converts small parameter changes into approximate outcome changes with clear sign and magnitude interpretations (\emph{marginal effects}).



\section{The univariate derivative: definition, meaning, and rules}

\subsection{Limit definition and economic interpretation}
\begin{definition}[Derivative]
Let \(y=f(x)\). The derivative of \(f\) at \(x\) is
\[
f'(x)=\lim_{h\to 0}\frac{f(x+h)-f(x)}{h},
\]
provided the limit exists. Economically, \(f'(x)\) is the \emph{marginal} effect of \(x\) on \(y\).
\end{definition}

\begin{example}[Marginal cost]
If total cost is \(C(Q)=a+bQ+cQ^2\) with \(a,b,c>0\), then the marginal cost is
\(
MC(Q)=C'(Q)=b+2cQ.
\)
MC rises with output when \(c>0\) (increasing marginal cost).
\end{example}

\subsection{Basic rules}
For differentiable functions \(u(x)\), \(v(x)\) and constant \(k\):
\[
\frac{\mathrm{d}}{\mathrm{d}x}\big[ku\big] = ku',\quad
\frac{\mathrm{d}}{\mathrm{d}x}\big[u\pm v\big] = u'\pm v',\quad
\frac{\mathrm{d}}{\mathrm{d}x}\big[uv\big]=u'v+uv',
\]
\[
\frac{\mathrm{d}}{\mathrm{d}x}\Big[\frac{u}{v}\Big]=\frac{u'v-uv'}{v^2},\qquad
\frac{\mathrm{d}}{\mathrm{d}x}f\!\big(g(x)\big)=f'\!\big(g(x)\big)\cdot g'(x) \quad(\text{chain rule}).
\]

\subsection{Higher-order derivatives, curvature, and Taylor}
\begin{definition}
The second derivative \(f''(x)\) measures how the slope changes. If \(f''(x)>0\) locally, \(f\) is \emph{convex}; if \(f''(x)<0\), \(f\) is \emph{concave}.
\end{definition}

A second-order Taylor approximation around \(x_0\) is
\[
f(x)\approx f(x_0)+f'(x_0)(x-x_0)+\tfrac12 f''(x_0)(x-x_0)^2,
\]
useful for local welfare and surplus approximations.

\subsection{Elasticities and log-derivatives}
The (point) elasticity of \(y\) with respect to \(x\) is
\[
\varepsilon_{y,x}=\frac{\mathrm{d}y}{\mathrm{d}x}\cdot\frac{x}{y}.
\]
If \(y=f(x)\) and both are positive, \(\dtotal{\ln y}{\ln x}=\varepsilon_{y,x}\). For power functions \(y=Ax^\alpha\), \(\varepsilon_{y,x}=\alpha\) everywhere.

\section{Partial and total derivatives}

\subsection{Partial derivatives}
For \(z=f(x_1,\dots,x_n)\), the partial derivative with respect to \(x_i\) is
\[
\pd{z}{x_i}=\lim_{h\to 0}\frac{f(x_1,\dots,x_i+h,\dots,x_n)-f(x_1,\dots,x_n)}{h},
\]
holding all other arguments fixed. Cross-partials \(\pd{^2 z}{x_i\,\partial x_j}\) (for \(i\neq j\)) capture interaction effects. If \(f\) is sufficiently smooth, \(\pd{^2 f}{x_i\,\partial x_j}=\pd{^2 f}{x_j\,\partial x_i}\) (Clairaut symmetry).

\subsection{Total derivative and total differential}
When \(x=x(p)\) depends on parameter \(p\),
\[
\frac{\mathrm{d}z}{\mathrm{d}p}=\sum_{i=1}^n \pd{f}{x_i}\frac{\mathrm{d}x_i}{\mathrm{d}p}+\pd{f}{p}.
\]
Infinitesimally,
\[
\mathrm{d}z=\sum_{i=1}^n \pd{f}{x_i}\,\mathrm{d}x_i+\pd{f}{p}\,\mathrm{d}p.
\]
These formulas are the workhorses of comparative statics when equilibria are described implicitly.

\section{Implicit differentiation and the implicit function theorem}

\subsection{Single-equation case}
Consider equilibrium defined by one equation \(F(x,p)=0\), where \(x\) is endogenous and \(p\) is a parameter.

\begin{proposition}[Chiang's implicit differentiation rule]
If \(F\) is continuously differentiable and \(F_x\neq 0\) at \((x^\ast,p)\), then locally there is a differentiable function \(x=x(p)\) with
\[
\boxed{\;\dtotal{x}{p}=-\frac{F_p}{F_x}\;}.
\]
\end{proposition}

\begin{example}[Linear supply--demand (parameter shift)]
Demand \(D(p,m)=a-bp+\alpha m\), supply \(S(p,w)=c+dp-\beta w\). Equilibrium: \(F(p;m,w)=D(p,m)-S(p,w)=0\).
Then \(F_p=-b-d<0\), \(F_m=\alpha>0\), \(F_w=-\beta<0\). Hence
\[
\dtotal{p^\ast}{m}=-\frac{F_m}{F_p}=\frac{\alpha}{b+d}>0,\qquad
\dtotal{p^\ast}{w}=-\frac{F_w}{F_p}=-\frac{-\beta}{b+d}=-\frac{\beta}{b+d}<0.
\]
Income raises the equilibrium price (normal good), higher input costs shift supply left and \emph{lower} the price here because \(w\) enters supply with a negative sign in \(F\); if one defines \(F=S-D\) instead, signs flip consistently. Always check which side of the equation your parameter enters.
\end{example}

\subsection{System case and Jacobians}
Let \(\ve x\in\R^n\) solve the \(n\)-equation system \(\ve F(\ve x,\ve p)=\ve 0\). Define the Jacobian with respect to endogenous variables
\[
J = \pd{\ve F}{\ve x}=
\begin{bmatrix}
\pd{F_1}{x_1} & \cdots & \pd{F_1}{x_n}\\
\vdots & \ddots & \vdots\\
\pd{F_n}{x_1} & \cdots & \pd{F_n}{x_n}
\end{bmatrix},\qquad
B=\pd{\ve F}{\ve p}.
\]
If \(J\) is nonsingular at the point of interest, then by the implicit function theorem
\[
\boxed{\;\dtotal{\ve x}{\ve p}=-J^{-1}B\; }.
\]

\paragraph{Cramer's rule for signs (2\(\times\)2 illustration).}
With two equations and one scalar parameter \(p\),
\[
\begin{aligned}
\pd{\ve F}{\ve x}
&=
\begin{bmatrix}
F_{1x_1} & F_{1x_2}\\
F_{2x_1} & F_{2x_2}
\end{bmatrix},\quad
B=
\begin{bmatrix}
F_{1p}\\
F_{2p}
\end{bmatrix},\quad
\Delta=\det(J)=F_{1x_1}F_{2x_2}-F_{1x_2}F_{2x_1}\neq 0,\\[4pt]
\dtotal{x_1}{p}
&=-\,\frac{\det\big(
\begin{smallmatrix}
\boxed{F_{1p}} & F_{1x_2}\\
\boxed{F_{2p}} & F_{2x_2}
\end{smallmatrix}\big)}{\Delta}
=-\,\frac{F_{1p}F_{2x_2}-F_{2p}F_{1x_2}}{\Delta},\qquad
\dtotal{x_2}{p}
=-\,\frac{\det\big(
\begin{smallmatrix}
F_{1x_1} & \boxed{F_{1p}}\\
F_{2x_1} & \boxed{F_{2p}}
\end{smallmatrix}\big)}{\Delta}
=-\,\frac{F_{1x_1}F_{2p}-F_{2x_1}F_{1p}}{\Delta}.
\end{aligned}
\]
Signs follow from the signs of cofactors and \(\Delta\).

\section{Worked applications}

\subsection{Tax incidence in linear supply--demand}
Let inverse demand \(p = \alpha - \beta Q\) (\(\beta>0\)) and inverse supply \(p = \gamma + \delta Q\) (\(\delta>0\)). Impose a per-unit \emph{consumer} tax \(t\), so consumers pay \(p_c=p+t\) while producers receive \(p\). Equilibrium satisfies
\[
\alpha - \beta Q = (\gamma + \delta Q) + t \quad\Rightarrow\quad Q^\ast(t)=\frac{\alpha-\gamma - t}{\beta+\delta}.
\]
Comparative statics:
\[
\dtotal{Q^\ast}{t}=-\frac{1}{\beta+\delta}<0,\qquad
\dtotal{p^\ast}{t}=\frac{\delta}{\beta+\delta}\in(0,1),\qquad
\dtotal{p_c^\ast}{t}=1-\frac{\delta}{\beta+\delta}=\frac{\beta}{\beta+\delta}\in(0,1).
\]
Price received by producers rises by a fraction \(\delta/(\beta+\delta)\) of the tax; the rest is borne by consumers. Incidence falls on the relatively inelastic side (larger slope in levels).

\subsection{Monopoly with linear demand and constant marginal cost}
Suppose inverse demand \(p=a-bQ\), \(a,b>0\), constant marginal cost \(c\in(0,a)\). Profit \(\pi(Q)=(a-bQ)Q-cQ-F\). FOC: \(\pi'(Q)=a-2bQ-c=0\Rightarrow Q^\ast=\dfrac{a-c}{2b}\), \(p^\ast=\dfrac{a+c}{2}\).
Comparative statics:
\[
\dtotal{Q^\ast}{c}=-\frac{1}{2b}<0,\quad
\dtotal{p^\ast}{c}=\frac12>0,\quad
\dtotal{Q^\ast}{a}=\frac{1}{2b}>0,\quad
\dtotal{p^\ast}{a}=\frac12>0.
\]
Interpretation: higher cost reduces output and raises price; stronger demand shifts increase both.

\subsection{Cournot duopoly (best-response system and Jacobian)}
Two firms choose quantities \(q_1,q_2\), price \(p=A-B(q_1+q_2)\), constant cost \(c_i\). FOCs:
\[
\begin{aligned}
F_1(q_1,q_2;c_1)&=A-B(q_1+q_2)-Bq_1-c_1=0 \ \Rightarrow\ A-c_1 -2Bq_1 -Bq_2=0,\\
F_2(q_1,q_2;c_2)&=A-c_2 -Bq_1 -2Bq_2=0.
\end{aligned}
\]
Jacobian \(J=\begin{bmatrix}-2B & -B\\ -B & -2B\end{bmatrix}\), \(\Delta=3B^2>0\).
Parameter matrix for \(c_1\): \(B_{(c_1)}=\big[\,-1,\ 0\,\big]^\top\).
By Cramer's rule,
\[
\dtotal{q_1^\ast}{c_1}=-\frac{\det\big(\begin{smallmatrix}\boxed{-1}&-B\\[1pt]\boxed{0}&-2B\end{smallmatrix}\big)}{\Delta}
=-\frac{2B}{3B^2}=-\frac{2}{3B}<0,\qquad
\dtotal{q_2^\ast}{c_1}=-\frac{\det\big(\begin{smallmatrix}-2B&\boxed{-1}\\[1pt]-B&\boxed{0}\end{smallmatrix}\big)}{\Delta}
=-\frac{-B}{3B^2}=\frac{1}{3B}>0.
\]
An increase in firm 1's cost reduces its own output and (strategic substitutability) raises its rival's.

\subsection{Cobb--Douglas demand and elasticities}
Utility \(u(x,y)=x^\alpha y^{1-\alpha}\) with \(0<\alpha<1\); prices \(p_x,p_y\); income \(I\).
Marshallian demand:
\[
x^\ast=\frac{\alpha I}{p_x},\qquad y^\ast=\frac{(1-\alpha)I}{p_y}.
\]
Comparative statics:
\[
\pd{x^\ast}{I}=\frac{\alpha}{p_x}>0,\quad \pd{x^\ast}{p_x}=-\frac{\alpha I}{p_x^2}<0,\quad
\pd{x^\ast}{p_y}=0.
\]
Elasticities: \(\varepsilon_{x,I}=1\) (homothetic), \(\varepsilon_{x,p_x}=-1\), and cross-price elasticity \(0\) (gross complements/substitutes determined by other preferences, here it is zero).

\subsection{A general single-equation nonlinearity}
Let equilibrium be \(F(x,p)=\ln x + x^2 - p=0\) with \(x>0\). Then
\[
F_x=\frac{1}{x}+2x>0,\qquad F_p=-1.
\]
Hence
\[
\dtotal{x^\ast}{p}=-\frac{F_p}{F_x}=\frac{1}{\frac{1}{x^\ast}+2x^\ast}=\frac{x^\ast}{1+2(x^\ast)^2}\in(0,\tfrac12),
\]
so \(x^\ast\) rises with \(p\) but at a diminishing rate.

\section{How-to guide for comparative statics}
\begin{tcolorbox}[sectionbox]
\textbf{Recipe.}
\begin{enumerate}[label=\arabic*.]
  \item \emph{Write} the equilibrium conditions \(F(\cdot)=0\) (or \(\ve F=\ve 0\)) with a clearly labeled endogenous vector \(\ve x\) and parameters \(\ve p\).
  \item \emph{Differentiate} implicitly: single equation \(\dtotal{x}{p}=-F_p/F_x\); system \(\dtotal{\ve x}{\ve p}=-J^{-1}B\).
  \item \emph{Sign} the derivatives using economically plausible slope conditions (e.g., \(D_p<0\), \(S_p>0\)) and stability (\(|J|\neq 0\), often \(|J|>0\)).
  \item \emph{Interpret} magnitudes with elasticities or linearization; check units.
  \item \emph{Sanity checks}: limit cases (perfectly inelastic/elastic), symmetry, and dimensions.
\end{enumerate}
\end{tcolorbox}

\section{Connections and extensions}
\paragraph{Envelope theorem (preview).} When an optimizer chooses \(x\) to maximize \(f(x,p)\), the effect of \(p\) on the optimal value is \(\pd{f}{p}\) evaluated at the optimum (no need to track \(\dtotal{x^\ast}{p}\)). This complements, but does not replace, comparative statics of \(x^\ast\).

\paragraph{Log-linear comparative statics.} If \(x^\ast(p)\!>\!0\), then \(\dtotal{\ln x^\ast}{p}=\dfrac{1}{x^\ast}\dtotal{x^\ast}{p}\) and \(\dtotal{\ln x^\ast}{\ln p}=\dfrac{p}{x^\ast}\dtotal{x^\ast}{p}\), useful for elasticity interpretations.

\section{Common mistakes}
\begin{itemize}
  \item Confusing \emph{partial} with \emph{total} derivatives: \(\pd{y}{p}\) holds other arguments fixed; \(\dtotal{y}{p}\) accounts for induced changes.
  \item Forgetting the sign of \(F_x\): in \(-F_p/F_x\), the denominator is the slope of the equilibrium condition with respect to the endogenous variable.
  \item Mixing definitions of \(F\): switching between \(D-S\) and \(S-D\) flips signs consistently; pick one and stick to it.
  \item Ignoring singular Jacobians: if \(|J|=0\), the local mapping may not exist; revisit model or check stability assumptions.
\end{itemize}

\section*{Summary}
\begin{itemize}
  \item Derivative \(f'(x)\): marginal effect; rules of differentiation; curvature via \(f''(x)\).
  \item Partial vs total derivatives; elasticities via log-derivatives.
  \item Implicit differentiation: single equation \(\dtotal{x}{p}=-F_p/F_x\).
  \item Systems: \(\dtotal{\ve x}{\ve p}=-J^{-1}B\); for \(2\times 2\), use Cramer's rule to sign effects.
  \item Applications: tax incidence, monopoly, Cournot, Cobb--Douglas.
\end{itemize}

\section*{Practice (with brief answers)}
\begin{enumerate}[label=\arabic*.]
\item Let \(F(x,p)=x^3-p=0\). Compute \(\dtotal{x}{p}\) at \(p>0\). \\
\emph{Answer:} \(F_x=3x^2\), \(F_p=-1\), so \(\dtotal{x}{p}=1/(3x^2)=1/(3p^{2/3})>0\).

\item IS--LM toy: \(Y=C_0+c(Y-T)+I_0-\alpha r+G\) and \(M/P=kY-hr\). Solve for \(\dtotal{Y}{G}\) sign using Jacobians with \(x_1=Y,x_2=r\).\\
\emph{Answer (sign):} Under \(0<c<1\), \(h,\alpha,k>0\), \(|J|>0\) and \(\dtotal{Y}{G}>0\).

\item Cournot with symmetric costs \(c\). Show \(\dtotal{q_1^\ast}{A}>0\), \(\dtotal{q_1^\ast}{B}<0\). \\
\emph{Answer:} From closed forms \(q_i^\ast=\dfrac{A-c}{3B}\), derivatives follow immediately.
\end{enumerate}

\bigskip
\hrule
\bigskip


\bigskip

\begin{thebibliography}{9}
\bibitem{CW}
Chiang, A.\,C. and Wainwright, K.\ (2005). \emph{Fundamental Methods of Mathematical Economics}, 4th ed., McGraw--Hill.

\bibitem{SB}
Simon, C.\,P. and Blume, L.\ (2010). \emph{Mathematics for Economists}. W.\,W. Norton.

\bibitem{SH}
Syds\ae ter, K. and Hammond, P.\ (2022). \emph{Essential Mathematics for Economic Analysis}, 6th ed., Pearson.
\end{thebibliography}

\end{document}
